% !TEX root = task-2-10-2.tex
% Editable task bank. Append new BankTask blocks at the end of the document.
\documentclass[12pt,a4paper]{article}
\usepackage{physics-textbook}
\usepackage{graphicx,needspace}
\usetikzlibrary{decorations.pathmorphing}
\newif\ifShowSolutions
\ifdefined\TutorEdition
  \ShowSolutionstrue
\else
  \ShowSolutionsfalse % Change to \ShowSolutionstrue to display solutions.
\fi
\newcounter{banktask}
\newsavebox{\BankTaskBox}
\newcommand{\BankTaskContent}[3]{%
  \noindent\textbf{\thebanktask.}\quad #2\par
  \ifShowSolutions
    \Needspace{50mm}% Keep the solution heading with its opening text or diagram.
    #3\par
    \ifdefined\TutorEdition
      \ifcsname TutorSource#1\endcsname
        \nopagebreak\smallskip{\small\csname TutorSource#1\endcsname\par}%
      \fi
    \fi
  \fi
}
\newcommand{\BankTask}[3]{%
  \refstepcounter{banktask}\par\medskip
  \sbox{\BankTaskBox}{\begin{minipage}{\linewidth}%
    \BankTaskContent{#1}{#2}{#3}\end{minipage}}%
  \ifdim\dimexpr\ht\BankTaskBox+\dp\BankTaskBox\relax>\dimexpr\textheight-20mm\relax
    \Needspace{35mm}\BankTaskContent{#1}{#2}{#3}%
  \else
    \Needspace{\dimexpr\ht\BankTaskBox+\dp\BankTaskBox\relax}%
    \noindent\usebox{\BankTaskBox}\par
  \fi
  \medskip
}
\begin{document}
\BookSubsection{2.10.2}{График плавления}


% Source 5745AB: https://ege.fipi.ru/bank/questions.php?search=1&proj=BA1F39653304A5B041B656915DC36B38&qid=5745AB
\BankTask{5745AB}{
\begin{center}
\begin{tikzpicture}[task graph,x=0.19cm,y=0.95cm]
\draw[task axis] (0,0)--(31,0) node[right] {$t,\,\text{мин.}$};
\draw[task axis] (0,0)--(0,4) node[left] {$T$};
\node[below left] at (0,0) {$0$};
\foreach \x in {5,10,15,20,25}{
  \node[below] at (\x,0) {$\x$};
}
\foreach \x/\y in {5/2.35,10/2.35,15/2.35,20/2.9,25/3.45}
  \draw[task projection] (\x,0)--(\x,\y);
\draw[task curve] (0,1.45)--(5,2.35)--(15,2.35)--(25,3.45);
\end{tikzpicture}
\end{center}
Для плавления куска льда при температуре его плавления требуется количество теплоты, равное \(3\,\text{кДж}\). Этот кусок льда внесли в тёплое помещение. Зависимость температуры льда от времени представлена на рисунке. Определите среднюю тепловую мощность, подводимую к куску льда в процессе плавления. (в $\text{Вт}$).
}{
\smallskip\noindent\textbf{Решение от GPT}\par
Плавлению соответствует горизонтальный участок графика от \(5\) до \(15\,\text{мин}\). Поэтому \(\Delta t=10\,\text{мин}=600\,\text{с}\), и средняя мощность равна
\[
P=\frac{Q}{\Delta t}=\frac{3000}{600}=5\,\text{Вт}.
\]
Ответ: \(5\,\text{Вт}\).
}

\end{document}
